% =====================================================================
%  Cross-Dataset External Validation of a Retinal Foundation Model
%  for Diabetic Retinopathy Grading: Performance, Explainability,
%  and Uncertainty
%
%  Manuscript draft. Written journal-agnostic (standard `article` class)
%  so it compiles anywhere with pdflatex. For final submission, port to
%  the target journal's class (e.g. MDPI mdpi.cls for Diagnostics /
%  Journal of Imaging) - see paper/SUBMISSION_CHECKLIST.md.
%
%  Build:  latexmk -pdf main.tex     (or: pdflatex; bibtex; pdflatex x2)
%
%  Placeholders marked  %% TODO(Phase 7/8/9)  are filled once those
%  phases complete. Numbers in the text are from Phases 2-6
%  (see ../PHASE*_REPORT.md); values marked \prov{} are PROVISIONAL
%  pending the RETFound full fine-tune re-run.
% =====================================================================
\documentclass[11pt,a4paper]{article}

\usepackage[utf8]{inputenc}
\usepackage[T1]{fontenc}
\usepackage[margin=1in]{geometry}
\usepackage{graphicx}
\usepackage{booktabs}
\usepackage{multirow}
\usepackage{amsmath}
\usepackage{siunitx}
\usepackage{authblk}
\usepackage{lineno}
\usepackage[hidelinks]{hyperref}
\usepackage[round]{natbib}
\usepackage{xcolor}

\newcommand{\todo}[1]{\textcolor{red}{\textbf{[TODO: #1]}}}
\newcommand{\prov}[1]{#1\textsuperscript{\textcolor{blue}{p}}} % provisional value

\title{Cross-Dataset External Validation of a Retinal Foundation Model
for Diabetic Retinopathy Grading: Performance, Explainability, and
Uncertainty}

\author[1]{\todo{Author One}}
\author[1,2]{\todo{Author Two}}
\author[1]{\todo{Corresponding Author}\thanks{Correspondence: \todo{email}}}
\affil[1]{\todo{Department, Institution, City, Country}}
\affil[2]{\todo{Second affiliation}}

\date{\today}

\begin{document}
\maketitle
\linenumbers

% ---------------------------------------------------------------------
\begin{abstract}
\noindent\textbf{Background:} Retinal foundation models such as RETFound are
promoted as improving generalisation of retinal image analysis, but
independent cross-dataset external validation for diabetic retinopathy (DR)
grading remains limited.
\textbf{Objective:} To test whether a fine-tuned retinal foundation model
generalises better than an ImageNet-pretrained convolutional neural network
(CNN) to independent external datasets for five-class DR severity grading, and
to characterise the explainability and predictive uncertainty of both models
under dataset shift.
\textbf{Methods:} An ImageNet-pretrained CNN baseline (EfficientNet-B3;
ResNet-50 as secondary) and RETFound (ViT-Large/16, masked-autoencoder
pretraining on colour fundus images) were fine-tuned on APTOS~2019
($n=\num{3662}$; stratified 70/10/20 train/validation/test split) using an
identical preprocessing pipeline, augmentation, loss, and early-stopping rule.
Both were then evaluated \emph{unchanged} on two external datasets, Messidor-2
($n=\num{1744}$) and IDRiD ($n=455$). The primary metric was the
quadratic-weighted kappa (QWK); secondary metrics were accuracy, macro-F1,
balanced accuracy, per-class sensitivity/specificity, and one-vs-rest AUC.
Explainability was assessed with gradient- and attention-based saliency against
a pre-specified clinical-plausibility rubric; predictive uncertainty was
quantified with Monte-Carlo dropout and evaluated for error detection and
selective prediction. Ninety-five percent confidence intervals were obtained by
bootstrap resampling; paired AUCs were compared with DeLong's test.
\textbf{Results:} On the internal APTOS test set the CNN baseline and RETFound
achieved comparable QWK (EfficientNet-B3 \num{0.876}; RETFound full fine-tune
\prov{\num{0.869}}). On external validation, \emph{all} models degraded
substantially, most severely on Messidor-2 (baseline QWK \num{0.499},
$\Delta=-\num{0.377}$; RETFound \prov{\num{0.368}},
$\Delta=\prov{-\num{0.501}}$) and less so on IDRiD (baseline \num{0.760},
$\Delta=-\num{0.116}$; RETFound \prov{\num{0.720}}, $\Delta=\prov{-\num{0.150}}$).
RETFound did not show a generalisation advantage over the ImageNet baseline on
either external dataset; external one-vs-rest AUC was comparable across models
($\approx$\numrange{0.75}{0.84}).
\todo{Phase 7/8/9 headline sentences: explainability plausibility rates;
uncertainty AUROC for error detection; selective-prediction gain; DeLong
$p$-values.}
\textbf{Conclusions:} In this retrospective, public-dataset study, fine-tuning a
retinal foundation model on APTOS~2019 did not improve cross-dataset
generalisation for DR grading relative to a standard ImageNet-pretrained CNN.
Large internal-to-external performance drops for every model underline that
external validation, not internal test performance, should govern claims of
clinical readiness.
\medskip

\noindent\textbf{Keywords:} diabetic retinopathy; deep learning; foundation
models; RETFound; external validation; domain shift; explainable AI;
uncertainty quantification; retinal imaging
\end{abstract}

% ---------------------------------------------------------------------
\section{Introduction}

Diabetic retinopathy (DR) is a leading cause of preventable vision loss, and
automated grading of colour fundus photographs has been studied extensively as a
route to scalable screening \citep{gulshan2016,ting2017}. Reported internal
performance is frequently high, but models often degrade when applied to images
from cameras, populations, or grading protocols not seen during training
\citep{beede2020,wang2020}. Rigorous external validation on independent datasets
is therefore the appropriate evidence standard, yet it is inconsistently
reported \citep{yang2024review,tripodai2024}.

Self-supervised ``foundation models'' pretrained on large unlabelled corpora
have been proposed to improve label efficiency and robustness. RETFound
\citep{zhou2023retfound} is a masked-autoencoder Vision Transformer pretrained
on approximately 1.6 million retinal images and reported to transfer well to
several ocular and systemic disease tasks. Whether this pretraining confers
better \emph{cross-dataset} generalisation for DR severity grading --
specifically relative to a strong ImageNet-pretrained convolutional baseline
under an identical fine-tuning and evaluation protocol -- has not been
independently established.

\paragraph{Research questions.}
\begin{enumerate}
  \item[\textbf{RQ1 (Performance).}] Does a fine-tuned retinal foundation model
    (RETFound) generalise better than an ImageNet-pretrained CNN to external
    datasets (Messidor-2, IDRiD) for five-class DR grading, measured primarily
    by the internal-to-external drop in quadratic-weighted kappa?
  \item[\textbf{RQ2 (Explainability).}] Are the image regions driving each
    model's predictions more clinically plausible (overlapping visible DR
    lesions rather than optic disc, vasculature, or image borders) for one
    model than the other, on external images?
  \item[\textbf{RQ3 (Uncertainty).}] Does predictive uncertainty (Monte-Carlo
    dropout) separate correctly from incorrectly classified external images,
    and can an uncertainty threshold support selective prediction that
    meaningfully raises accuracy on accepted cases?
\end{enumerate}

\paragraph{Contribution.} We report a fully reproducible external validation in
which one preprocessing pipeline, one training protocol, and one evaluation
protocol are applied identically to an ImageNet CNN baseline and to RETFound. We
provide the master results table across one internal and two external datasets,
an explainability analysis against an a-priori clinical rubric, an uncertainty
analysis with selective-prediction curves, and full code and split definitions.
Our central finding is negative: under matched conditions, RETFound did not
outperform the CNN baseline on external DR grading.

% ---------------------------------------------------------------------
\section{Materials and Methods}

This study followed the TRIPOD+AI reporting guideline \citep{tripodai2024}; the
completed checklist is provided as Supplementary Material
(\texttt{paper/TRIPOD-AI\_checklist.md}).

\subsection{Datasets}

Three public datasets were used (Table~\ref{tab:datasets}). Full provenance
(source URL, version, access date, licence, citation) is recorded in the Data
Provenance Log (\texttt{DATA\_PROVENANCE.md}).

\begin{itemize}
  \item \textbf{APTOS 2019 Blindness Detection} \citep{aptos2019}: \num{3662}
    macula-centred colour fundus images from Aravind Eye Hospital, India, graded
    0--4 on the International Clinical DR severity scale. Used \emph{only} for
    model development (training, validation, internal test).
  \item \textbf{Messidor-2} \citep{decenciere2014,abramoff2013}: \num{1744}
    images with an adjudicated reference standard \citep{krause2018}. External
    test only.
  \item \textbf{IDRiD (Disease Grading subset)} \citep{porwal2018,porwal2020}:
    a 455-image subset (of the 516-image official release; the local mirror was
    incomplete, reported as a limitation) with ICDR 0--4 grades. External test
    only.
\end{itemize}

Messidor-2 and IDRiD were not inspected or used for any modelling or
hyperparameter decision until the single external evaluation
(Section~\ref{sec:external}).

\begin{table}[t]
\centering
\caption{Dataset characteristics. Grade share is the percentage of gradable
images at each ICDR severity level.}
\label{tab:datasets}
\small
\begin{tabular}{lrlccccc}
\toprule
Dataset & $n$ & Role & \multicolumn{5}{c}{Grade share (\%): 0 / 1 / 2 / 3 / 4} \\
\midrule
APTOS 2019 & \num{3662} & Development   & 49.3 & 10.1 & 27.3 & 5.3 & 8.1 \\
Messidor-2 & \num{1744} & External test & 58.3 & 15.5 & 19.9 & 4.3 & 2.0 \\
IDRiD      & 455        & External test & 28.4 & 4.8  & 34.3 & 18.5 & 14.1 \\
\bottomrule
\end{tabular}
\end{table}

\subsection{Preprocessing}

A single preprocessing script was applied identically to all three datasets:
(i) conversion to RGB; (ii) tolerance-based cropping of the black frame around
the circular fundus; (iii) zero-padding to a centred square (aspect ratio
preserved); (iv) resizing to $224\times224$ (RETFound's native input). Channel
normalisation used ImageNet statistics ($\mu=(0.485,0.456,0.406)$,
$\sigma=(0.229,0.224,0.225)$), consistent with RETFound's downstream
configuration. Contrast-limited adaptive histogram equalisation was implemented
as an optional step and was \emph{not} applied in the primary analysis. All
\num{5861} preprocessed images passed a stratified visual audit for
label--image alignment.

\subsection{Development split}

APTOS 2019 was partitioned 70\%/10\%/20\% into training, internal validation,
and internal test, stratified by DR grade, with a fixed random seed (42); the
exact image-ID lists are provided
(\texttt{data\_splits/aptos\_\{train,val,test\}.csv}).

\subsection{Models and fine-tuning}

\paragraph{Baseline CNN.} EfficientNet-B3 (primary) and ResNet-50 (secondary),
ImageNet-pretrained (\texttt{torchvision}), with the classifier head replaced by
a five-way head. Full fine-tuning (all layers unfrozen), AdamW
($\text{lr}=\num{1e-4}$, weight decay \num{1e-4}), cosine annealing, batch
size 32, mixed precision, up to 30 epochs with early stopping on internal-
validation QWK (patience 7). The loss was cross-entropy weighted by inverse
class frequency. Augmentation: horizontal/vertical flips, $\pm15^\circ$
rotation, mild brightness/contrast jitter.

\paragraph{RETFound.} The ViT-Large/16 masked-autoencoder checkpoint pretrained
on colour fundus images \citep{zhou2023retfound} was loaded (encoder weights
only; pretrained final LayerNorm reused as the pooling norm; fresh five-way
head). Two variants were fine-tuned, forming the principal ablation:
(a) \emph{linear probing} -- backbone frozen, head trained (base
$\text{lr}=\num{1e-3}$, up to 40 epochs); (b) \emph{full fine-tuning} -- all
layers unfrozen with layer-wise learning-rate decay (0.65), drop-path 0.2, label
smoothing 0.1, base $\text{lr}=\num{1e-4}$, batch size 16, up to 30 epochs.
Data, normalisation, class weighting, metric set, and early-stopping rule were
identical to the baseline. All hyperparameters are logged
(\texttt{outputs/phase\{4,5\}\_*/\,*\_hyperparams.json}).

\subsection{Evaluation metrics}

Given the marked class imbalance in every dataset, the primary metric was the
quadratic-weighted kappa (QWK), which credits ordinal proximity. Secondary
metrics: overall accuracy, macro-averaged F1, balanced accuracy, per-class
sensitivity and specificity, and one-vs-rest AUC (per class and macro-averaged).
The internal-to-external \emph{drop} in each metric is reported per model per
external dataset.

\subsection{External validation}
\label{sec:external}

Each trained model was run once, unmodified and without any further training or
threshold tuning, on Messidor-2 and on IDRiD, and re-scored on the APTOS
internal test set through the identical inference code path. Per-image outputs
(true label, predicted label, class probabilities, predictive entropy) were
retained for all model $\times$ dataset combinations.

\subsection{Explainability}
%% TODO(Phase 7): finalise once maps are produced.
For the CNN, Grad-CAM++ maps were generated for the final convolutional block;
for the ViT, attention-rollout maps from the last block were used. A matched
sample of correctly and incorrectly classified external images across DR grades
was rated against a pre-specified rubric: a map was \emph{plausible} if its peak
overlapped a visible microaneurysm, haemorrhage, or exudate, and
\emph{implausible} if concentrated on the optic disc, major vessels, or image
border with no visible lesion. \todo{Second-rater agreement (Cohen's $\kappa$)
or single-rater limitation statement.}

\subsection{Uncertainty quantification}
%% TODO(Phase 8): finalise once MC-dropout run completes.
Monte-Carlo dropout was enabled at inference ($T=30$ stochastic forward passes);
predictive entropy of the mean softmax was the uncertainty score. We tested
whether misclassified external images had higher uncertainty than correctly
classified ones (Mann--Whitney $U$), reported the AUROC of entropy for error
detection, and plotted selective-prediction curves (accuracy versus coverage as
the most-uncertain predictions are referred for human review).

\subsection{Statistical analysis}
%% TODO(Phase 9): fill p-values / CI half-widths.
Ninety-five percent confidence intervals for all headline metrics were obtained
by bootstrap resampling of the test set with replacement (\num{2000}
iterations, percentile method). Paired AUCs for the baseline versus RETFound on
the same dataset were compared with DeLong's test. The internal-to-external drop
was reported with its own bootstrap CI for each model so that the
\emph{difference} in drop between models could be assessed. $\alpha = 0.05$,
two-sided. Analyses used \texttt{scikit-learn}, \texttt{scipy}, and
\texttt{statsmodels} in Python 3.14.

\subsection{Reproducibility}

All code (preprocessing, split generation, training, evaluation, statistics),
the exact split image-ID lists, per-image prediction files, and an environment
specification are available (Section~\ref{sec:availability}). Raw datasets are
not redistributed; each is linked to its source with licence terms in
\texttt{DATA\_PROVENANCE.md}.

% ---------------------------------------------------------------------
\section{Results}

\subsection{Internal performance}

On the APTOS internal test set ($n=733$), the CNN baseline and RETFound achieved
broadly comparable performance (Table~\ref{tab:master}). EfficientNet-B3 reached
QWK \num{0.876} (accuracy \num{0.804}); RETFound full fine-tuning reached QWK
\prov{\num{0.869}} (accuracy \prov{\num{0.769}}) and linear probing QWK
\num{0.847}. RETFound showed flatter per-class sensitivity, with higher recall
on the rare proliferative grade than the CNNs.

\subsection{External validation and performance drop}

\begin{table}[t]
\centering
\caption{Master results: each model evaluated unchanged on the internal test and
two external datasets. QWK = quadratic-weighted kappa; Bal-Acc = balanced
accuracy; AUC = one-vs-rest macro AUC. Values marked \textsuperscript{p} are
PROVISIONAL pending completion of the RETFound full fine-tune re-run.
95\% CIs to be added (Phase 9).}
\label{tab:master}
\small
\begin{tabular}{llrrrrrr}
\toprule
Model & Dataset & $n$ & Accuracy & QWK & Macro-F1 & Bal-Acc & AUC \\
\midrule
\multirow{3}{*}{EfficientNet-B3 (baseline)}
 & APTOS-internal & 733  & 0.804 & 0.876 & 0.644 & 0.641 & 0.928 \\
 & Messidor-2     & 1744 & 0.626 & 0.499 & 0.360 & 0.353 & 0.775 \\
 & IDRiD          & 455  & 0.497 & 0.760 & 0.441 & 0.497 & 0.830 \\
\midrule
\multirow{3}{*}{ResNet-50 (baseline)}
 & APTOS-internal & 733  & 0.789 & 0.875 & 0.616 & 0.630 & 0.924 \\
 & Messidor-2     & 1744 & 0.614 & 0.435 & 0.354 & 0.339 & 0.754 \\
 & IDRiD          & 455  & 0.448 & 0.707 & 0.403 & 0.477 & 0.841 \\
\midrule
\multirow{3}{*}{RETFound (linear probe)}
 & APTOS-internal & 733  & 0.759 & 0.847 & 0.607 & 0.633 & 0.908 \\
 & Messidor-2     & 1744 & 0.589 & 0.394 & 0.347 & 0.423 & 0.758 \\
 & IDRiD          & 455  & 0.336 & 0.626 & 0.325 & 0.395 & 0.778 \\
\midrule
\multirow{3}{*}{RETFound (full fine-tune)\textsuperscript{p}}
 & APTOS-internal & 733  & \prov{0.769} & \prov{0.869} & \prov{0.616} & \prov{0.629} & \prov{0.915} \\
 & Messidor-2     & 1744 & \prov{0.602} & \prov{0.368} & \prov{0.362} & \prov{0.346} & \prov{0.758} \\
 & IDRiD          & 455  & \prov{0.420} & \prov{0.720} & \prov{0.423} & \prov{0.480} & \prov{0.821} \\
\bottomrule
\end{tabular}
\end{table}

\begin{table}[t]
\centering
\caption{Internal-to-external drop in QWK. Negative values indicate degradation
on the external dataset.}
\label{tab:drop}
\small
\begin{tabular}{lrrr}
\toprule
Model & QWK internal & QWK Messidor-2 ($\Delta$) & QWK IDRiD ($\Delta$) \\
\midrule
EfficientNet-B3 (baseline) & 0.876 & 0.499 ($-0.377$) & 0.760 ($-0.116$) \\
ResNet-50 (baseline)       & 0.875 & 0.435 ($-0.440$) & 0.707 ($-0.167$) \\
RETFound (linear probe)    & 0.847 & 0.394 ($-0.453$) & 0.626 ($-0.221$) \\
RETFound (full fine-tune)\textsuperscript{p} & \prov{0.869} & \prov{0.368} (\prov{$-0.501$}) & \prov{0.720} (\prov{$-0.150$}) \\
\bottomrule
\end{tabular}
\end{table}

Every model degraded on both external datasets (Tables~\ref{tab:master},
\ref{tab:drop}; Figure~\ref{fig:forest}). The degradation was largest on
Messidor-2, where QWK fell from $\approx$\num{0.87} internally to
\numrange{0.37}{0.50}. On Messidor-2, all models reverted toward majority-class
(grade~0) prediction: grade-1 (mild) sensitivity was $\le \num{0.02}$ for every
model, and grade-0 specificity fell to \numrange{0.18}{0.32}. On IDRiD the QWK
drop was smaller (\numrange{-0.12}{-0.22}) but accuracy fell more
(\numrange{-0.31}{-0.35}), consistent with IDRiD's pathology-enriched class mix
(Table~\ref{tab:datasets}).

\paragraph{RQ1.} RETFound did not generalise better than the ImageNet CNN
baseline on either external dataset. Its external QWK was lower than the
EfficientNet-B3 baseline on both Messidor-2 and IDRiD, and its full-fine-tune
Messidor-2 drop was the largest of any model (\prov{$-0.501$} vs $-0.377$).
External one-vs-rest AUC was comparable across all models
($\approx$\numrange{0.75}{0.84}). \todo{DeLong $p$-values (Phase 9): state
whether the AUC differences are non-significant.} RETFound did show isolated
per-class advantages under shift -- highest severe-grade sensitivity on IDRiD
and highest proliferative-grade sensitivity on Messidor-2 -- and, under linear
probing, the highest balanced accuracy on Messidor-2, indicating a more even
error distribution across grades.

\begin{figure}[t]
\centering
%% TODO(Phase 10): \includegraphics[width=0.9\linewidth]{../outputs/figN_forest.pdf}
\fbox{\parbox{0.8\linewidth}{\centering\vspace{2cm}
Forest / bar plot of QWK, accuracy and AUC for each model across the internal
and two external datasets, with 95\% CI error bars (Phase 10, from
\texttt{outputs/phase6/master\_results.csv} + Phase 9 bootstrap CIs).%
\vspace{2cm}}}
\caption{Model performance across datasets with 95\% confidence intervals.}
\label{fig:forest}
\end{figure}

\subsection{Explainability}
%% TODO(Phase 7)
\todo{Plausibility rates (\% of maps overlapping a visible lesion) for baseline
vs RETFound, correctly vs incorrectly classified, per external dataset; 4--6
representative panels as Figure 3; rater-agreement or single-rater caveat.}

\subsection{Uncertainty}
%% TODO(Phase 8)
\todo{Mann--Whitney $U$ result and AUROC of predictive entropy for detecting
external misclassifications, per model per dataset; selective-prediction curve
(Figure 4) and the accuracy gain / error capture at, e.g., 90\% coverage.}

% ---------------------------------------------------------------------
\section{Discussion}

\paragraph{RQ1 -- Performance.} Under a matched preprocessing, fine-tuning, and
evaluation protocol, fine-tuning RETFound on APTOS~2019 did not yield better
cross-dataset generalisation for five-class DR grading than an
ImageNet-pretrained CNN. Both RETFound variants had lower external QWK than the
EfficientNet-B3 baseline, and the direction of the result was consistent across
linear probing (which ran to a clean early stop) and full fine-tuning. This does
not contradict RETFound's reported strengths on other tasks, but it tempers the
expectation that retinal-domain pretraining automatically confers robustness to
the specific shifts represented by Messidor-2 and IDRiD (different cameras,
acquisition, populations, and -- for Messidor-2 -- an adjudicated reference
standard).

\paragraph{Nature of the shift.} The Messidor-2 collapse is an ordinal-structure
failure: models predict almost exclusively grade~0 or grade~2, so accuracy is
partly preserved by the high grade-0 prevalence while QWK falls sharply. The
IDRiD pattern is different -- a base-rate shift toward more advanced disease --
so accuracy falls more than QWK. Reporting balanced accuracy alongside QWK
(Table~\ref{tab:master}) separates these effects, and we recommend this practice
for external-validation studies generally.

\paragraph{RQ2 / RQ3.} \todo{One paragraph each once Phases 7--8 complete:
whether either model's saliency was more lesion-focused under shift, and whether
uncertainty was actionable for triage (e.g. ``referring the most-uncertain
10\% of external predictions for human review would capture X\% of the errors'').}

\paragraph{Comparison with the literature.} \todo{Contrast the observed
internal-to-external QWK drops with values reported (or not reported) in the
Phase~1 literature table, e.g. \citep{yang2024review} and primary
external-validation reports for DR grading.}

\paragraph{Limitations.}
(i) The RETFound full fine-tune reported here was interrupted before its planned
schedule completed because of an infrastructure I/O bottleneck; a completed run
(in progress) is expected to update these values, although the concordant
linear-probe result makes a reversal of RQ1 unlikely. The final manuscript will
use the completed run.
(ii) All datasets are retrospective and public; no prospective or
clinical-deployment claim is made.
(iii) The IDRiD mirror used contained 455 of 516 Disease-Grading images.
(iv) External sample sizes limit precision for rare grades (grade-1 sensitivity
is not estimable on Messidor-2 for any model); wide CIs are expected (Phase 9).
(v) \todo{Single-rater plausibility assessment, if a second rater is not
obtained (Phase 7).}
(vi) A single foundation model (RETFound MAE, colour-fundus weights) and a
single downstream task were studied.

\paragraph{Clinical and practical implications.} The magnitude of the
internal-to-external drop -- for every model, foundation or not -- reinforces
that internal test performance is a poor guide to behaviour on new data, and
that external validation and, where possible, an uncertainty-aware referral
mechanism should precede any screening deployment.

% ---------------------------------------------------------------------
\section{Conclusions}

In this retrospective external-validation study, a fine-tuned retinal foundation
model did not outperform a standard ImageNet-pretrained CNN for cross-dataset
diabetic retinopathy grading. All models -- foundation and baseline alike --
lost substantial ordinal agreement on independent data, most severely on
Messidor-2. These results argue for external validation as the primary evidence
standard and motivate uncertainty-aware deployment safeguards.

% ---------------------------------------------------------------------
\section*{Declarations}

\subsection*{Data and Code Availability}
\label{sec:availability}
Analysis code, split image-ID lists, per-image prediction files, and the
environment specification are available at \todo{repository URL / DOI}. The three
datasets are third-party and are not redistributed: APTOS~2019 via Kaggle
\citep{aptos2019}; Messidor-2 via ADCIS \citep{decenciere2014} with adjudicated
grades from \citep{krause2018}; IDRiD via IEEE DataPort \citep{porwal2018}. See
\texttt{DATA\_PROVENANCE.md} for URLs, versions, access dates, and licences.

\subsection*{Ethics}
This study used previously published, fully de-identified, publicly available
datasets collected under the original studies' approvals and consent
procedures; no new human-subjects data were collected. \todo{Insert the target
journal's required wording for secondary use of public data; confirm whether an
institutional ethics waiver statement is required.}

\subsection*{Funding} \todo{Funding sources or ``This research received no
external funding.''}

\subsection*{Conflicts of Interest} \todo{Declare or ``The authors declare no
conflict of interest.''}

\subsection*{Author Contributions} \todo{CRediT roles.}

\subsection*{Use of AI tools} \todo{Disclose any use of generative-AI assistance
in code or drafting, per the target journal's policy.}

% ---------------------------------------------------------------------
\bibliographystyle{plainnat}
\bibliography{references}

\end{document}
